\chapter{Financial reports across regimes}\label{ch:financialreport}

\section{Introduction}

Three regimes require the periodic reports of listed companies to be tagged: ESEF in the European Union, the Inline XBRL mandate of the U.S. Securities and Exchange Commission, and the XBRL that Japan's Financial Services Agency has required of a securities report since 2008. Modules \emph{esef}, \emph{sec} and \emph{edinet} of chapters \ref{ch:esef} to \ref{ch:edinet} each bring one of them into the lake as its publisher gives it out, in the shape that publisher gives it. Module \emph{combined} of the package \emph{financialreport} reduces them to the same columns and writes one table of facts and one of filings, so that a question put to all three regimes is put to one table rather than to three of different shapes.

It holds no data of its own and collects nothing. There is no Bronze layer, nothing is downloaded, and both tables are derived entirely from tables that are themselves rebuilt, so both are rebuilt in full on every run, with no watermark to keep and nothing to resume from. The dependency runs one way: the source modules do not import this one and do not know it exists, their tables stay authoritative and complete, and a lake that has collected only ESEF filings loses nothing by building nothing here. Which sources a run combines is decided by which of them have tables to read. A lake with one source gets a combined table of that source, which is the same question answered over the population there is, and a lake with none is told so rather than given empty tables.

The pieces fit together because the source modules spell the columns they share in the same way. A fact has, in every regime, a filing, a concept, the first and last day of its period, a unit and a value; this module adds which source the row came from and the year of the period the filing reports on, by which, after the source, the table of facts is partitioned. What one regime does not give out is left empty rather than filled: the company files of the Commission carry no dimensions and no text, so those columns are empty for every SEC row, and an empty dimension there means a figure for the entity as a whole.

\section{The entity}

The tables do not join the regimes by entity, because they cannot. ESEF prescribes the LEI, a filing with the Commission states a central index key, and an EDINET filing an EDINET code, and no body publishes a mapping between the first two. GLEIF records the register an entity is registered in, and for some tens of thousands of records that register is EDGAR, but nearly all of those are funds named by a series identifier; the record of an operating company points at its state register instead, NVIDIA's at Delaware by a file number that is not its central index key. Each row therefore carries an LEI, a central index key or an EDINET code, exactly one of them is filled, and none is derived from another. The honest shape of the table says that the join is not available rather than filling a column with a guess.

Japan is the partial exception described in section \ref{sec:edinet-lei}. The corporate number that the agency's list of filers carries is a check digit followed by the number of the commercial register, which GLEIF records against an LEI issued to a Japanese company, so the join is exact where it exists. The combined tables carry the corporate number beside the EDINET code for that reason. It reaches about one listed Japanese filer in seven, the rest having no LEI at all, so it is enough to put a Japanese filer beside a European one where both exist and not enough to key anything by.

\section{The framework}

Which accounting framework a figure was prepared under is a column of its own, and it is not the regime the filing was made under. For ESEF and SEC rows it is read from the taxonomy the concept is tagged against: \verb|IFRS| for \verb|ifrs-full|, \verb|US-GAAP| for \verb|us-gaap|. That is what makes usable the population in which the regimes genuinely meet: a foreign private issuer reporting under IFRS files with the Commission and tags against the IFRS taxonomy, and belongs beside the European reports rather than beside the American ones it was filed with.

For EDINET rows the prefix does not settle it. A Japanese group that has adopted IFRS tags its statements against Japan's own IFRS taxonomy and its front matter against the same disclosure taxonomy every other filer uses, and its parent company's own statements are Japan GAAP in the same report: Fujitsu states revenue of 3,502,971 million for the group under the one and 1,692,618 million for the company under the other. The framework of a Japanese fact is therefore taken from its taxonomy where the taxonomy is decisive -- IFRS for the Japanese IFRS taxonomy, \verb|JP-GAAP| for the taxonomy of the financial statements ordinance -- and from the standard the filing states otherwise. A filing to which no accounting standard applies, such as a large-holding report or a fund's report, has no framework rather than the conversion's mark for not applicable.

Two properties of the Japanese rows have to be read before they are used, and both belong to the regime. A report that consolidates states every figure twice, the group's without a dimension and the parent company's under a member for the non-consolidated basis, and both are in the table because both are in the source; a question that does not say which it wants is asking for two sets of accounts at once. And Japan GAAP is not comparable with the other two line by line: its profit and loss account has an ordinary result between the operating result and the result before tax, and extraordinary items below it, which an IFRS statement does not have. The intermediate lines are not comparable; the bottom line still is.

One kind of fact is not carried, and it is the only thing the module drops. A fact of the EDINET table whose period could not be resolved has no end date, and a fact in no period can be compared with nothing, which is what this table is for. Such facts stay in the EDINET table of facts, where a check of that module counts them.

\section{The catalogue of measures}

Whether \verb|us-gaap:Revenues| and \verb|ifrs-full:Revenue| are the same figure is a judgement about accounting rather than about storage. The module makes it once, in one place where it can be read and disagreed with, and as data rather than code: a catalogue shipped with the package that defines eight measures -- assets, liabilities, equity, equity attributable to the parent, profit or loss, profit or loss attributable to the parent, revenue, and cash and cash equivalents -- in 38 rows, of which 29 carry a caveat. A row names a measure, a framework, a taxonomy prefix and an element, a priority among the concepts of one framework, the reason for the choice, where the frameworks do not actually agree, and who decided. Nothing in it is authoritative. There is no official correspondence between IFRS and US GAAP, and what is shipped is a starting point that a deployment should review and extend with a catalogue of the same shape in its own package, read beside this one.

The catalogue is used in two ways, and the difference between them is the point of the design. A function resolves a measure at query time: for each filing and period it takes the best concept the filing states, by the priorities, excludes facts reported under a dimension, and warns where the concepts it resolved through carry caveats. This answers a question about a population and is what a benchmark should use. A column of the table of facts, \emph{measure}, instead names what a fact is, beside the concept its filer tagged, which is never altered. It is held to a stricter rule: only the first concept of each framework carries the name, and every other fact carries none. A measure earns the column only once its concepts are agreed to mean the same figure in every framework the catalogue defines it for, because a query filtering on it is making the claim that the rows it returns are comparable. All eight measures pass that rule today, but a measure added to the catalogue is not materialised by being added; it has to be named as such, so that writing a row and claiming a figure is comparable stay two separate acts.

The stricter rule has costs, and they are not symmetric between regimes. The first concepts of equity and of profit or loss are the totals including non-controlling interests. Many American registrants tag only the figures that exclude them, which for a registrant without minorities carry the same number, and the element does not say whether there are any; so those facts carry no name. Every Japanese filer, by contrast, tags its profit or loss, and the population under the name is close to all of them. The two samples are selected by different rules, and a comparison of their medians is not a comparison of the two regimes. The measure of profit attributable to the parent leans the opposite way, so between them the two measures cover the population well, but neither does alone.

One gap is structural rather than a matter of what a filer tagged. Japan GAAP has no element for the equity attributable to the parent, which is a sum of two lines of the balance sheet rather than a line, so no Japanese filing under Japan GAAP can carry that measure and a query filtering on it excludes Japan without saying so. The sum is not computed to close the gap: taking the nearest element instead is short by the valuation and translation adjustments, 9,191 million of a correct 63,326 million on one report, and adding the components would be the module inventing a figure, which no module of the package does. For the same reason a missing total is never derived from the balance sheet identity, and a count of the rows carrying a measure is a count of what was stated and named, never a population.

The same figure is not always in the same place. A Japanese balance sheet has no line for cash and cash equivalents; it presents cash and deposits, including time deposits of more than three months, which are cash equivalents under none of the three frameworks. The figure that means what \verb|ifrs-full:CashAndCashEquivalents| means is the closing balance of the Japanese cash flow statement, and that carries the name.

\section{Data pipeline}

There is no landing zone and no Bronze layer. The Silver layer holds the two tables. The table of facts has one row per fact of a filing, partitioned by source and then by the year of the period the filing reports on, since most questions are scoped to one regime and a question about a figure to a year after that. The table of filings has one row per filing of every source, with the entity keys, the name, the framework, the member state for an ESEF filing, the form for a filing with the Commission or the agency, the period end and how many facts the filing contributed. Like the tables of its sources it lists filings that were never read, so a count of what was read filters on that number. There is no Gold layer: the module stops at resolution. It says which concepts a measure is taken to be and labels facts accordingly, but does not define ratios, aggregate across entities or build a metrics layer, which belong to whoever is asking the question. The sector factsheet of chapter \ref{ch:esef} is one such reader, reaching from ESEF to issuers filing with the Commission or on EDINET through these tables.

\section{Software architecture}

Module \emph{combined} has three parts. The part \emph{common} declares the sources, the layout and the schemas, and describes what each source covers; \emph{mapping} reads the catalogue of measures and resolves measures over any table of facts spelled as the sources spell them -- including the table of a single source, so that an ESEF-only lake can resolve a measure without this module's tables; and \emph{silver} reduces each source to the common columns and writes the two tables in one call,
\begin{quote}
\verb|combined.update_silver_financialreport(spark, "paths.yaml")|
\end{quote}
which a scheduler makes after the source pipelines have run. The configuration file needs a section for the module naming its Silver layer and nothing else, together with the sections of the sources it reads.

Two report templates of the module read the tables as Parquet and render with Quarto to PDF. A \emph{coverage report} says what each regime's part of the lake holds -- filings listed and read, issuers, years and frameworks, the forms read, and which of the eight measures the annual reports of each regime and framework state. A \emph{comparison report} sets the companies of the regimes beside one another on three ratios of figures of the same report, net margin, equity ratio and return on equity, so that no currency is converted, and prints the caveats of the catalogue on the measures it uses. Both leave out what is not a company: a Japanese investment fund files its annual report under the same document type as a company, and in a lake that read a few weeks of the agency's indexes most annual reports are funds', 456 of 605 in the one the reports were written against; the comparison leaves out banks and insurers as well. Both read a Japanese filer without subsidiaries from the accounts it states under \verb|NonConsolidatedMember|, and count a measure only under the framework of the filing, since a Japanese group reporting under IFRS states its parent company's own accounts under Japan GAAP in the same report. The benchmark of the ESEF sector factsheet, which reaches the Japanese filings through the combined tables, applies the same three rules and sets Japanese banks, insurers and securities firms apart by the industry elements they tag.
